% Standalone proof dossier for the perimeter/excess consumption cluster.
\documentclass[../main.tex]{subfiles}
\begin{document}

% === SOLUTION HEADER =========================================================
% Certification is not recorded here: research/program/ledger.yaml owns the
% proofs[] mode and review path, and the review's front matter owns identities.
%
% ledger-node : prop:intro-audit; prop:trivial-excess; lem:perimeter-martingale; lem:excess-identity; lem:inf-martingales
% refines : prop:intro-audit (modules/kls/20-eldan-statements.tex);
% author : /root/kls_bootstrap_author
% date : 2026-08-27
% ============================================================================
\section*{Solution: the perimeter martingale and the excess-consumption audit}
\label{sec:sol-kls-excess-audit}

\paragraph{Scope and regularity.}
Let $(\mu_t)_{t\geq0}$ be Eldan's stochastic-localization posterior started from a
log-concave probability measure $\mu$ on $\mathbb R^n$.  We use
\[
  F_t(x)=\frac{\dd\mu_t}{\dd\mu}(x),\qquad
  \dd F_t(x)=F_t(x)(x-a_t)\mathbin\cdot\dd W_t,
  \qquad a_t=\int x\,\dd\mu_t(x).
\]
For a Borel set $E$ and $r>0$, put
\[
  E_r=\{x\in\mathbb R^n:\operatorname{dist}(x,E)<r\},
\]
and use the lower outer Minkowski convention of the manuscript,
\begin{equation}
  P_t(E)=\mu_t^+(E)
  :=\liminf_{r\downarrow0}\frac{\mu_t(E_r)-\mu_t(E)}r.
  \label{eq:sol-minkowski-perimeter}
\end{equation}
We also write
\[
  p_t=\mu_t(E),\qquad
  e_t(E)=P_t(E)-I_{\mu_t}(p_t).
\]
The general fixed-cut result below assumes $P_0(E)<\infty$ and uses only bounded-test mass
martingales.  The stronger martingale identity is stated separately in the compact-support
smooth regularity class: the initial density is smooth and compactly supported and $E$ has a
$C^2$ boundary with a tubular neighborhood over the support of the density, so the one-sided
tube formula identifies
\eqref{eq:sol-minkowski-perimeter} with weighted surface area.  No rough-set approximation or
identification of inequivalent perimeter conventions is used.

\subsection*{1. Perimeter under localization}

\begin{lemma}[Perimeter martingale; Lemma~\ref{lem:perimeter-martingale}]
\label{lem:sol-perimeter-martingale}
For every Borel set $E$ with $P_0(E)<\infty$, the lower Minkowski perimeter process is an
integrable nonnegative supermartingale:
\begin{equation}
  \mathbb E[P_t(E)\mid\mathcal F_s]\leq P_s(E),\qquad 0\leq s\leq t.
  \label{eq:sol-perimeter-supermartingale}
\end{equation}
In the compact-support smooth regularity class specified above, $P_t(E)$ is instead a true
martingale and
\begin{equation}
  \dd P_t(E)
  =\left(\int_{\partial^*E}(x-a_t)\,\dd\sigma_t(x)\right)\mathbin\cdot\dd W_t,
  \label{eq:sol-perimeter-sde}
\end{equation}
where $\sigma_t$ is the weighted surface measure of $\mu_t$ on $\partial^*E$.  Thus its drift
vanishes identically and $\mathbb E P_t(E)=P_0(E)$.
\end{lemma}

\begin{proof}
We first prove the general lower-Minkowski assertion directly.  For $q\in\mathbb Q_{>0}$ define
\[
  X_{q,t}=\frac{\mu_t(E_q)-\mu_t(E)}q
         =\frac{\mu_t(E_q\setminus E)}q.
\]
For each fixed $q$, this is $q^{-1}$ times the mass of a fixed Borel set.  The bounded-test
localization identity \eqref{eq:loc-martingale} therefore makes
$(X_{q,t})_{t\geq0}$ a bounded, nonnegative true
martingale.  For $n\geq1$ let
\[
  Y_{n,t}=\inf\{X_{q,t}:q\in\mathbb Q,\ 0<q<1/n\}.
\]
The infimum is over a nonempty countable fixed family.  For every admissible $q$ and
$0\leq s\leq t$,
\[
  \mathbb E[Y_{n,t}\mid\mathcal F_s]
  \leq\mathbb E[X_{q,t}\mid\mathcal F_s]=X_{q,s}.
\]
Because there are only countably many $q$, the inequalities may be realized on one common
full-probability set.  Infimizing there gives
\begin{equation}
  \mathbb E[Y_{n,t}\mid\mathcal F_s]\leq Y_{n,s}.
  \label{eq:sol-boundary-layer-supermartingale}
\end{equation}
Also $0\leq Y_{n,t}\leq X_{q_n,t}$ for any one admissible rational $q_n$, so $Y_{n,t}$ is
integrable.

For any probability measure $\nu$, the map $r\mapsto\nu(E_r\setminus E)$ is continuous from
below: if $q_k\uparrow r$, then $E_{q_k}\setminus E\uparrow E_r\setminus E$.  Approximating
every real $r>0$ from below by rationals consequently shows that the infimum over rational
$q\in(0,1/n)$ equals the infimum over all real $r\in(0,1/n)$.  Hence, pathwise,
\begin{equation}
  Y_{n,t}\uparrow
  \liminf_{r\downarrow0}\frac{\mu_t(E_r)-\mu_t(E)}r=P_t(E).
  \label{eq:sol-boundary-layer-limit}
\end{equation}
Ordinary monotone convergence and \eqref{eq:sol-boundary-layer-supermartingale} with $s=0$
give
\[
  \mathbb E P_t(E)
  =\lim_n\mathbb E Y_{n,t}
  \leq\lim_nY_{n,0}=P_0(E)<\infty.
\]
Thus $P_t(E)$ is integrable.  Conditional monotone convergence in
\eqref{eq:sol-boundary-layer-supermartingale} now yields
\[
  \mathbb E[P_t(E)\mid\mathcal F_s]
  =\lim_n\mathbb E[Y_{n,t}\mid\mathcal F_s]
  \leq\lim_nY_{n,s}=P_s(E),
\]
which proves \eqref{eq:sol-perimeter-supermartingale} for arbitrary support and without a
rough-set approximation.

It remains to prove the stronger statement in the compact-support smooth class, using the same
Minkowski convention.  Let $w_0$ be the smooth compactly supported density of $\mu$ and put
$\dd\sigma_0=w_0\,\dd\mathcal H^{n-1}|_{\partial^*E}$.  The one-sided tubular-neighborhood
formula for a smooth boundary and the continuous density $w_t=F_tw_0$ gives, for every fixed
$t$ and sample path,
\begin{equation}
  P_t(E)=\int_{\partial^*E}F_t(x)\,\dd\sigma_0(x),
  \qquad
  \dd\sigma_t(x)=F_t(x)\,\dd\sigma_0(x).
  \label{eq:sol-perimeter-representation}
\end{equation}
Indeed, in outward normal coordinates the Jacobian is $1+O(r)$ and
$w_t(x+r\mathbf n_x)\to w_t(x)$, so division of the boundary-layer integral by its thickness
and dominated convergence gives the displayed surface integral.  Thus
\eqref{eq:sol-perimeter-representation} is precisely the lower Minkowski perimeter
\eqref{eq:sol-minkowski-perimeter}, not a change of convention.  Its initial surface measure is
finite.

Let $K$ contain the support of $\mu$ and set $D=\operatorname{diam}K$.  Both $x$ and $a_u$
belong to the convex hull of $K$ for $\sigma_0$-almost every $x$, so $|x-a_u|\leq D$.  For fixed
$x$, let $\rho_m\uparrow\infty$ localize the stochastic integral in the pointwise SDE and apply It\^o's
formula to $F_{u\wedge\rho_m}(x)^2$.  Taking expectations gives
\[
  \mathbb E F_{u\wedge\rho_m}(x)^2
  \leq 1+D^2\int_0^u\mathbb E F_{v\wedge\rho_m}(x)^2\,\dd v.
\]
Gronwall and then Fatou as $m\uparrow\infty$ give
\begin{equation}
  \mathbb E F_u(x)^2\leq \exp(D^2u).
  \label{eq:sol-density-L2}
\end{equation}
Consequently, for every finite $T$,
\begin{align*}
  \mathbb E\int_0^T\int_{\partial^*E}
    |F_u(x)(x-a_u)|^2\,\dd\sigma_0(x)\,\dd u
  &\leq D^2\sigma_0(\partial^*E)\int_0^T e^{D^2u}\,\dd u<\infty.
\end{align*}
Since $\sigma_0$ is finite, Cauchy--Schwarz also gives
\[
  \mathbb E\int_0^T
  \left|\int_{\partial^*E}F_u(x)(x-a_u)\,\dd\sigma_0(x)\right|^2\dd u<\infty.
\]
These are pre-interchange stochastic-Fubini bounds.  Stochastic Fubini applied to the pointwise
density SDE therefore proves \eqref{eq:sol-perimeter-sde};
$F_t\,\dd\sigma_0=\dd\sigma_t$ gives its displayed coefficient.

Finally, the same bound (or Novikov directly) makes every fixed-$x$ process $F_t(x)$ a true
martingale on bounded intervals.  Conditional Tonelli in
\eqref{eq:sol-perimeter-representation} gives
\[
  \mathbb E[P_t(E)\mid\mathcal F_s]
  =\int_{\partial^*E}\mathbb E[F_t(x)\mid\mathcal F_s]\,\dd\sigma_0(x)
  =P_s(E).
\]
Thus the smooth compact-support lower Minkowski perimeter is a true martingale with the stated
SDE.
\end{proof}

\subsection*{2. The unconditional excess bound and the exact identity}

\begin{proposition}[Unconditional integrated excess; Proposition~\ref{prop:trivial-excess}]
\label{prop:sol-trivial-excess}
Let $\mu$ be isotropic and log-concave, let $E$ be Borel with
$\mu(E)=1/2$ and $P_0(E)<\infty$, and put $e_0=P_0(E)-I_\mu(1/2)$.  For every $T>0$ and every
stopping time $\tau$,
\begin{equation}
  \mathbb E\int_0^{T\wedge\tau}e_t(E)\,\dd t\leq(1+e_0)T.
  \label{eq:sol-trivial-excess}
\end{equation}
\end{proposition}

\begin{proof}
The profile is nonnegative, so $0\leq e_t(E)\leq P_t(E)$ pathwise.  Consequently, Tonelli and
\eqref{eq:sol-perimeter-supermartingale} give
\[
  \mathbb E\int_0^{T\wedge\tau}e_t(E)\,\dd t
  =\int_0^T\mathbb E[\mathbf 1_{\{t<\tau\}}e_t(E)]\,\dd t
  \leq\int_0^T\mathbb E P_t(E)\,\dd t
  \leq T P_0(E).
\]
This calculation uses neither optional stopping nor independence of $\tau$; only the
deterministic-time supermartingale bound is used.

To bound $P_0(E)$, take a coordinate marginal of $\mu$.  It is a one-dimensional log-concave
law $\nu$ of variance one.  If $f$ is its density and $m$ a median, the coordinate halfspace
$H=\{x:x_1\leq m\}$ has $\mu(H)=1/2$ and $\mu^+(H)=f(m)$.  Bobkov's
one-dimensional isoperimetric identity and variance estimate
\cite[Proposition~4.1 and (4.2)]{Bobkov1999LogConcave} give
$\operatorname{Is}(\nu)=2f(m)$ and
$\operatorname{Is}(\nu)^2\leq2/\operatorname{Var}_\nu(X)$, hence
$I_\mu(1/2)\leq f(m)\leq1/\sqrt2<1$.  Therefore
$P_0(E)=I_\mu(1/2)+e_0\leq1+e_0$, proving
\eqref{eq:sol-trivial-excess}.
\end{proof}

\begin{lemma}[Exact excess identity; Lemma~\ref{lem:excess-identity}]
\label{lem:sol-excess-identity}
Under the compact-support smooth regularity hypotheses of
Lemma~\ref{lem:sol-perimeter-martingale}, for every finite $t\geq0$,
\begin{equation}
  \mathbb E e_t(E)
  =e_0(E)+I_\mu(p_0)-\mathbb E I_{\mu_t}(p_t).
  \label{eq:sol-excess-identity}
\end{equation}
\end{lemma}

\begin{proof}
All quantities are nonnegative and integrable in the compact-support class.  By definition,
\[
  \mathbb E e_t(E)=\mathbb E P_t(E)-\mathbb E I_{\mu_t}(p_t).
\]
Lemma~\ref{lem:sol-perimeter-martingale} gives
$\mathbb E P_t(E)=P_0(E)$, while
$P_0(E)=I_\mu(p_0)+e_0(E)$.  Substitution proves
\eqref{eq:sol-excess-identity}.  Compact support is essential for the equality used here: the
noncompact argument above supplies only $\mathbb E P_t(E)\leq P_0(E)$.
\end{proof}

\subsection*{3. What a countable infimum of fixed-cut perimeter supermartingales implies}

\begin{lemma}[Fixed competitor families; Lemma~\ref{lem:inf-martingales}]
\label{lem:sol-inf-martingales}
Let $\mathfrak S$ be a fixed nonempty countable family of Borel sets of finite initial lower
Minkowski perimeter and define
$J_t=\inf_{S\in\mathfrak S}P_t(S)$.  Then $J$ is a supermartingale.  The same conclusion holds
when a prescribed uncountable family has a fixed nonempty countable determining subfamily
$\mathfrak S_0$ satisfying
$\inf_{S\in\mathfrak S}P_t(S)=\inf_{S\in\mathfrak S_0}P_t(S)$ almost surely at every time under
consideration; in that case $J$ means the latter countable infimum.  On
$\{t<\tau_\eta\}$ one has the pointwise comparison
\begin{equation}
  I_{\mu_t}(p_t)\geq J_t(\eta),
  \qquad
  J_t(\eta)=\inf\{P_t(S):\mu_t(S)\in[1/2-\eta,1/2+\eta]\}.
  \label{eq:sol-window-infimum}
\end{equation}
The family in the last infimum depends on $(t,\omega)$; the first assertion therefore does not
say that $J_t(\eta)$ or $I_{\mu_t}(p_t)$ is a supermartingale.
\end{lemma}

\begin{proof}
Fix $0\leq s\leq t$.  For every $S\in\mathfrak S$, $J_t\leq P_t(S)$, hence by monotonicity of
conditional expectation and the noncompact part of
Lemma~\ref{lem:sol-perimeter-martingale},
\[
  \mathbb E[J_t\mid\mathcal F_s]
  \leq\mathbb E[P_t(S)\mid\mathcal F_s]
  \leq P_s(S).
\]
Taking the infimum in $S$ on the right gives
$\mathbb E[J_t\mid\mathcal F_s]\leq J_s$.  This argument needs only that each fixed-cut
perimeter is a supermartingale; in the compact-support class the second inequality happens to
be an equality.  Moreover, choosing one $S_\star\in\mathfrak S$ gives
$0\leq J_t\leq P_t(S_\star)$ and
$\mathbb E P_t(S_\star)\leq P_0(S_\star)<\infty$, so $J_t$ is integrable.
The determining-subfamily variant is exactly the same countable argument.  No conclusion is
claimed for a raw uncountable pointwise infimum merely assumed measurable.

On $\{t<\tau_\eta\}$, $p_t\in[1/2-\eta,1/2+\eta]$.  Thus the exact-mass competitor class
$\{S:\mu_t(S)=p_t\}$ is contained in the windowed class in
\eqref{eq:sol-window-infimum}.  The infimum over the smaller class is at least the infimum over
the larger one, proving the pointwise comparison.  But eligibility for the windowed class is
the random condition $\mu_t(S)\in[1/2-\eta,1/2+\eta]$; it is not a fixed family to which the
conditional-expectation argument can be applied.
\end{proof}

\subsection*{4. Consumption audit}

\begin{proposition}[Consumption audit; Proposition~\ref{prop:intro-audit}]
\label{prop:sol-intro-audit}
The following three assertions hold.
\begin{enumerate}
  \item For every isotropic log-concave $\mu$, every Borel $E$ with $\mu(E)=1/2$ and
  $P_0(E)<\infty$, every $T>0$, and every stopping time $\tau$,
  \[
    \mathbb E\int_0^{T\wedge\tau}e_t(E)\,\dd t\leq(1+e_0)T.
  \]
  \item Replace the covariance weight in the stable Stein-trace estimate by $1$, keeping its
  fixed tight window and absorption margin.  Then this unweighted estimate, without any
  separate excess-propagation hypothesis, implies KLS.
  \item There are log-concave pairs $(\nu,E)$ with $\nu(E)=1/2$, $r=D=0$,
  $e(E)\to0$, and $\mathcal S_\nu(E)/s\to\infty$.  Hence the unweighted estimate cannot be
  proved by a universal one-time-slice inequality of the asserted form.
\end{enumerate}
\end{proposition}

\begin{proof}
The first assertion is Proposition~\ref{prop:sol-trivial-excess}.

For the second, spell out exactly what is consumed.  Fix a universal
$\eta\in(0,1/6]$ and constants $C_0,C_1,C_2$, $\beta\geq0$ satisfying
$2\beta+64\eta^2<1$.  The unweighted stable Stein-trace estimate is
\begin{equation}
\begin{split}
  \mathbb E\int_0^{T\wedge\tau_\eta}
       \frac{\mathcal S_{\mu_t}(E)}{s_t}\,\dd t
  \leq{}&C_0T+C_1\mathbb E\int_0^{T\wedge\tau_\eta}r_t\,\dd t
       +\beta\mathbb E\int_0^{T\wedge\tau_\eta}D_t\,\dd t\\
  &+C_2\mathbb E\int_0^{T\wedge\tau_\eta}e_t(E)\,\dd t .
\end{split}
\label{eq:sol-unweighted-stein}
\end{equation}
The exact tight-window Stein/source conversion
(Lemma~\ref{lem:stein-vs-source}) is
\[
  \mathbb E\int_0^{T\wedge\tau_\eta}S_t\,\dd t
  \leq2\mathbb E\int_0^{T\wedge\tau_\eta}
          \frac{\mathcal S_{\mu_t}(E)}{s_t}\,\dd t
       +64\eta^2\mathbb E\int_0^{T\wedge\tau_\eta}D_t\,\dd t.
\]
For the near-Cheeger cuts used in the contradiction argument, take $e_0\leq1$.  Applying the
first assertion to \eqref{eq:sol-unweighted-stein} gives
\[
\begin{split}
  \mathbb E\int_0^{T\wedge\tau_\eta}S_t\,\dd t
  \leq{}&(2C_0+4C_2)T
       +2C_1\mathbb E\int_0^{T\wedge\tau_\eta}r_t\,\dd t\\
       &+(2\beta+64\eta^2)
         \mathbb E\int_0^{T\wedge\tau_\eta}D_t\,\dd t.
\end{split}
\label{eq:sol-absorptive-carleson}
\]
This is precisely an absorptive tight-window two-color Carleson estimate, with damping
coefficient strictly below one.  Corollary~\ref{cor:tight-window-consumption} then gives a
positive universal boundary lower bound for every balanced near-Cheeger cut to which
\eqref{eq:sol-unweighted-stein} applies.

For completeness, if KLS failed, there would be isotropic log-concave $\mu_k$ with
$h_{\mu_k}\to0$.  The balanced characterization from Lemma~\ref{lem:half},
$h_\nu=2I_\nu(1/2)$, permits choices of sets $E_k$ of mass $1/2$ with
$P_0(E_k)\leq I_{\mu_k}(1/2)+o(1)$.  Thus $P_0(E_k)\to0$ and $e_0(E_k)\leq1$ for all large
$k$, whereas \eqref{eq:sol-absorptive-carleson} and tight-window consumption give a uniform
positive lower bound for the same perimeters, a contradiction.  No weighted excess estimate
was used.

The third assertion is exactly the certified anisotropic two-tail construction of
Proposition~\ref{prop:two-tail}: it has $r=D=0$, vanishing absolute excess, and divergent
$\mathcal S/s$.  A slice-wise version of \eqref{eq:sol-unweighted-stein} would have a uniformly
bounded right-hand side on that family and a divergent left-hand side.  This disproves that
slice-wise route while making no claim against nonlocal-in-time proofs or the covariance-weighted
formulation.
\end{proof}

\paragraph{Obstructions respected.}
No node in this dossier has a ledger \texttt{bounded\_by} edge.  Nevertheless,
Lemma~\ref{lem:sol-excess-identity} and Lemma~\ref{lem:sol-inf-martingales} explicitly preserve
the content of \texttt{obs:circularity}: the exact identity is restricted to compact support,
and no supermartingale property is inferred for the random time-dependent balanced-profile
infimum.  Proposition~\ref{prop:sol-intro-audit} uses the separate two-tail obstruction only
to refute a slice-wise unweighted estimate; it does not promote numerical or directional
evidence to proof.

\end{document}
